\documentclass{assets/utmreport}

\usepackage{booktabs,amssymb,float,graphicx,adjustbox,array}

\hypersetup{hidelinks,hypertexnames=false,pdfauthor={Lucian-Adrian Gavril}}

\selectlanguage{english}

\utmlogo{assets/utm-logo.jpeg}

\institution{Technical University of Moldova}
\faculty{Faculty of Computers, Informatics, and Microelectronics}
\department{Department of Software Engineering}
\studyprogramme{Software Engineering Study Programme}
\documenttype{Report on laboratory work no. 9}
\reporttitle{Operating Systems}
\reportsubtitle{Watching System Calls with strace}
\student{Lucian-Adrian Gavril}{FAF-241}
\supervisor{Reițman Patricia}{univ. assist.}
\cityyear{Chișinău}{2026}

\newcommand{\shot}[2]{\begin{figure}[H]\centering\adjustbox{scale=0.6,max width=\linewidth,max height=0.78\textheight}{\includegraphics{assets/#1}}\caption{#2}\end{figure}}

\begin{document}
\sloppy
\makeutmreporttitlepage
{\small\tableofcontents}
\clearpage

\utmintroduction
\textbf{Purpose of the work.} To watch the boundary between a program and the kernel. \texttt{strace} prints every system call a process makes, so the work reads the calls of a real command (\texttt{ls -l}), writes a program that uses \texttt{open}, \texttt{write}, \texttt{read} and \texttt{close} directly and checks that each C call becomes exactly one system call, and counts one kind of call across several programs to see what decides how many times a program enters the kernel.

\textbf{Environment.} Debian GNU/Linux 13, Linux kernel 6.12, GCC 14, strace 6.13. All commands ran in an ordinary user shell. Every terminal figure is the real output of the commands shown in it; the raw traces are kept in \texttt{evidence/}.

\chapter{The system calls of a real command}

\section{strace on ls -l}
\texttt{strace -f -e trace=openat,read,write,close ls -l} prints the four selected calls beside the normal output (Figure~1.1). The first part is the dynamic loader and the C library: \texttt{openat} on \texttt{ld.so.cache} and on each shared library, each library followed by one \texttt{read} of its ELF header and a \texttt{close}. Then \texttt{ls} opens the locale files, \texttt{/etc/nsswitch.conf}, \texttt{/etc/passwd} and \texttt{/etc/group} to turn the numeric owner and group into names for the long listing. The listing itself leaves the process in a single call at the end.
\shot{fig1_1_strace_ls.png}{strace -f -e trace=openat,read,write,close ls -l}

\section{How many calls}
\texttt{strace -c} gives the exact count instead of an estimate (Figure~1.2). With the same filter, one \texttt{ls -l} makes \textbf{87} calls: 39 \texttt{openat}, 33 \texttt{close}, 14 \texttt{read} and 1 \texttt{write}. Without any filter it makes \textbf{214} system calls in total, most of them \texttt{mmap}, \texttt{fstat}, \texttt{statx} and \texttt{getxattr} (one \texttt{statx} and the extended-attribute lookups per listed file). Part of the \texttt{openat} calls fail with \texttt{ENOENT}: these are locale paths that \texttt{ls} tries before the one that exists.
\shot{fig1_2_count_ls.png}{Exact call counts for ls -l, filtered and unfiltered}

\section*{Observe 1}
\textbf{Which system call does ls use to print to the screen? To which descriptor does it write?}
\texttt{ls} prints with the \texttt{write} system call, on descriptor \textbf{1}, standard output: \texttt{write(1, "total 16\textbackslash n...", 250) = 250} in Figure~1.1. The whole listing goes out in that one call, because the C library collects the lines in the \texttt{stdout} buffer and flushes them together when \texttt{ls} exits.

\chapter{A program that calls the kernel directly}

\section{Build and run}
\texttt{syscalls.c} is the program from the brief, with the return values of \texttt{open}, \texttt{write} and \texttt{read} also checked, so a failure prints \texttt{perror} instead of continuing with a bad descriptor. It builds with \texttt{-Wall -Wextra} and no warnings, and prints exactly the expected three lines (Figure~2.1). \texttt{note.txt} holds the 14 bytes \texttt{hello, kernel\textbackslash n}.
\shot{fig2_1_build_run.png}{Build with -Wall -Wextra and run of syscalls}

\section{The same program seen by strace}
\texttt{strace -e trace=openat,write,read,close ./syscalls} shows the six calls of \texttt{main} after the loader lines (Figure~2.2). Each C call is one system call with the same arguments: \texttt{open(...)} becomes
\begin{center}\texttt{openat(AT\_FDCWD, "note.txt", O\_WRONLY|O\_CREAT|O\_TRUNC, 0644) = 3}\end{center}
because glibc implements \texttt{open} with \texttt{openat} relative to the current directory; \texttt{write(3, "hello, kernel\textbackslash n", 14) = 14}; \texttt{close(3) = 0}; the second \texttt{openat} gets 3 again, because 3 was freed by the \texttt{close}; \texttt{read(3, ..., 63) = 14} asks for 63 bytes (\texttt{sizeof(buf) - 1}) and receives the 14 that exist.

The three \texttt{printf} lines do not appear as three writes. They appear at the end as one \texttt{write(1, ..., 80)}: when \texttt{stdout} is not a terminal (here strace's output is redirected), the C library buffers the whole output and writes it once at \texttt{exit}.
\shot{fig2_2_strace_syscalls.png}{strace of syscalls: the program's own calls after the loader}

\section*{Observe 2}
\textbf{The descriptor returned by the first openat is 3, not 0, 1 or 2. Why are those three already taken?}
Every process starts with three descriptors already open, inherited from the shell through \texttt{fork} and kept through \texttt{exec}: 0 is standard input, 1 standard output, 2 standard error. In a terminal all three point to the terminal device. The kernel always returns the \emph{lowest free} descriptor, and the lowest free one is 3.

Figure~2.3 proves both halves of this. \texttt{ls -l /proc/self/fd} shows 0, 1 and 2 pointing to \texttt{/dev/pts/2}. Then \texttt{./syscalls 0<\&-} starts the program with descriptor 0 closed, and \texttt{open} returns \textbf{0}: the number is not reserved for anything, it was simply in use.
\shot{fig2_3_fds.png}{Descriptors 0, 1, 2 of a process, and open returning 0 when stdin is closed}

\chapter{Counting one kind of call}

\section{The whole profile of syscalls}
\texttt{strace -c ./syscalls} counts every call the program makes (Figure~3.1): \textbf{41} in total. Only six belong to \texttt{main}; the rest is the work of starting a C program: \texttt{execve}, 3 \texttt{brk}, 8 \texttt{mmap}, 3 \texttt{mprotect}, \texttt{openat} and \texttt{read} of \texttt{libc.so.6}, \texttt{arch\_prctl}, \texttt{set\_tid\_address}, \texttt{rseq} and others. The counts match the program exactly: 2 \texttt{write} (one to \texttt{note.txt}, one flush of \texttt{stdout}), 2 \texttt{read} (\texttt{libc} header and \texttt{note.txt}), 4 \texttt{openat} and 4 \texttt{close}.
\shot{fig3_1_count_c.png}{strace -c ./syscalls: 41 system calls}

\section{write across programs}
\texttt{strace -c -e trace=write} counts only \texttt{write} (Figure~3.2). \texttt{ls -l /} makes 1; \texttt{seq 1 100000} makes 72 to send 588895 bytes, that is one call per full 8~KiB buffer (about 8192 bytes each); a shell loop of five \texttt{echo} commands makes 5, one per \texttt{echo}, because \texttt{echo} is a shell built-in that writes its line immediately. The number of \texttt{write} calls depends on how the program buffers, not on how many lines it prints.
\shot{fig3_2_count_write.png}{Number of write calls for ls, seq and a shell loop}

\section{Buffering decides the count}
The same three-line \texttt{printf} program behaves differently depending on where \texttt{stdout} goes (Figure~3.3). Into a pipe the C library uses full buffering: \textbf{one} \texttt{write} of 80 bytes at exit. On a terminal (run under \texttt{script}, so \texttt{stdout} is a tty) it uses line buffering: \textbf{three} writes of 27, 20 and 33 bytes, one per newline. This is why the strace of Figure~2.2 shows one write to descriptor 1, while a person running \texttt{./syscalls} in a terminal sees each line as soon as it is printed.
\shot{fig3_3_buffering.png}{The same program: one write into a pipe, three on a terminal}

\section*{Answers in brief}
\begin{tabular}{@{}p{0.34\linewidth}p{0.58\linewidth}@{}}
\toprule
Question & Answer \\
\midrule
How many calls does \texttt{ls -l} make? & 87 openat/read/write/close; 214 system calls in total \\
How does \texttt{ls} print? & One \texttt{write} on descriptor 1 (250 bytes) \\
Why is the first descriptor 3? & 0, 1, 2 are stdin, stdout, stderr, inherited open; the kernel gives the lowest free one \\
Calls of \texttt{./syscalls} & 41 total, 6 from \texttt{main} \\
\bottomrule
\end{tabular}

\utmconclusions
A system call is visible and countable. strace showed that every \texttt{open}, \texttt{write}, \texttt{read} and \texttt{close} in my program became exactly one call into the kernel with the same arguments, and that the library translated \texttt{open} into \texttt{openat}. It also showed how much happens before \texttt{main}: 35 of the 41 calls of \texttt{syscalls} are the loader mapping the C library.

Descriptors are small integers that index the process's table of open files, and the kernel always hands out the lowest free one. The fact that my file got 3 is not a rule of the API: when I closed standard input before starting the program, the same \texttt{open} returned 0.

The count of \texttt{write} calls is decided by buffering. \texttt{ls} sends its whole listing in one write, \texttt{seq} sends 588895 bytes in 72 writes, and my three \texttt{printf} lines cost one write into a pipe but three on a terminal. A system call is expensive compared with a function call, so the C library batches output, and strace is the tool that makes that visible.

\end{document}
